% !TEX root = ../reporte_final.tex
\section{Metodología}

El experimento sigue el flujo canónico de minería de datos KDD \citep{fayyad1996kdd}, instanciado como una sucesión de siete etapas explícitas que coinciden con las fases solicitadas en el enunciado del proyecto.

\subsection{Pipeline general}

\begin{enumerate}
  \item \textbf{Definición del problema.} Recuperar la estructura de las tres especies de pingüinos a partir de variables morfométricas, sin supervisión.
  \item \textbf{Adquisición de los datos.} Descarga del archivo \texttt{penguins.csv} del repositorio oficial \texttt{allisonhorst/palmerpenguins} (CC-0). En el notebook, si el archivo local no existe se obtiene automáticamente vía \texttt{urllib.request}, garantizando ejecución en Google Colab sin pasos manuales.
  \item \textbf{Análisis exploratorio.} Conteo por especie, \texttt{describe()}, recuento de \texttt{NaN}, matriz de correlación y \emph{pairplot} discriminado por clase.
  \item \textbf{Preprocesamiento.} Eliminación de filas con \texttt{NaN} en variables numéricas y en \texttt{species}; estandarización de las cuatro variables numéricas con \texttt{StandardScaler}.
  \item \textbf{Diseño de experimentos.} Aplicación de K-Means ($k=3$, $n_{init}=10$), DBSCAN ($\varepsilon=0.6$, $\text{minPts}=5$) y Agglomerative ($k=3$, linkage=Ward); semilla \texttt{RANDOM\_STATE} = 42 en todos los pasos estocásticos.
  \item \textbf{Evaluación.} Cálculo de Silhouette Score (excluyendo los puntos de ruido de DBSCAN), Adjusted Rand Index contra la etiqueta verdadera, matriz de confusión del mejor modelo y visualización en el plano PCA.
  \item \textbf{Conclusiones.} Discusión de los resultados, limitaciones y trabajo futuro.
\end{enumerate}

\subsection{Pseudocódigo del experimento}

\begin{lstlisting}[language=Python,caption={Pipeline esencial de minería de datos.}]
RANDOM_STATE = 42
df = read_csv("data/penguins.csv")
df_clean = df.dropna(subset=NUM_COLS + ["species"]).reset_index(drop=True)
X = StandardScaler().fit_transform(df_clean[NUM_COLS])
y_true = LabelEncoder().fit_transform(df_clean["species"])

models = {
  "K-Means":       KMeans(n_clusters=3, n_init=10, random_state=RANDOM_STATE),
  "DBSCAN":        DBSCAN(eps=0.6, min_samples=5),
  "Agglomerative": AgglomerativeClustering(n_clusters=3, linkage="ward"),
}

for name, m in models.items():
    labels = m.fit_predict(X)
    sil  = silhouette_score(X[labels != -1], labels[labels != -1])
    ari  = adjusted_rand_score(y_true, labels)
    log(name, sil, ari)
\end{lstlisting}

\subsection{Variables del dataset}

La tabla~\ref{tab:variables} resume las variables empleadas en el modelado. Únicamente las cuatro numéricas alimentan los algoritmos; \texttt{species} se reserva como verdad de campo para la métrica externa.

\begin{table}[H]
\centering
\caption{Variables del dataset Palmer Penguins consideradas en el experimento.}
\label{tab:variables}
\begin{tabularx}{\textwidth}{lXl}
\toprule
\textbf{Variable} & \textbf{Descripción} & \textbf{Tipo}\\
\midrule
\texttt{bill\_length\_mm}     & Longitud del pico (culmen) en milímetros          & Numérica continua\\
\texttt{bill\_depth\_mm}      & Profundidad del pico en milímetros                & Numérica continua\\
\texttt{flipper\_length\_mm}  & Longitud de la aleta en milímetros                & Numérica continua\\
\texttt{body\_mass\_g}        & Masa corporal en gramos                            & Numérica continua\\
\texttt{species}              & \emph{Adelie}, \emph{Chinstrap} o \emph{Gentoo}   & Categórica (verdad)\\
\bottomrule
\end{tabularx}
\end{table}
